\section{Introduction}

When a government reserves most of its low-value procurement for small firms and makes opting out of the reservation a scrutinized administrative act, what is the mechanism behind the resulting price increase, and when does that mechanism make the policy defensible? The reduced-form answer on prices---somewhere between ten and fifteen percent on winning prices---has been in hand for a decade \citep{marion2007, krasnokutskaya2011}. What price regressions cannot recover is whether the markup runs mainly through bidding behavior inside the restricted pool or through changes in who enters, and they cannot translate the markup into allocative waste plus tax distortion. This paper builds the structure needed for those splits and uses them to compare bidder exclusion with a price-preference alternative.

The setting is S\~ao Paulo state's BEC platform, the centralized electronic reverse auction through which the state procures standardized goods. On March 1, 2018, a legal opinion from the state prosecutor's office (PGE-SP) reversed a four-year administrative exception and extended the SME-only regime to medical supplies (Group 65). The decision was internal to PGE-SP and defended on isonomy grounds; there is no evidence it responded to prices, competition, or supplier complaints within Group 65. An earlier version of this paper used the same cutoff to estimate, by difference-in-differences, that winning prices rose 10--11 percent on Group-65 items relative to 76 never-treated product groups. Those numbers are the reduced form. What they do not reveal is whether the markup ran mostly through bidders or mostly through bidding, and whether the entry response mitigated or amplified the welfare loss.

To answer that, I specify an asymmetric independent-private-values model of the Preg\~ao format. Bidders draw costs from type-specific distributions $F_c^{\text{SME}}$ and $F_c^{\neg}$ and compete under a descending-clock English-reverse auction in which losers exit at cost and the last remaining firm supplies the item. Losers' drop-out prices therefore point-identify $F_c^k$ under the standard \citet{haile2003} interpretation, and winners are treated as left-censored at the second-lowest exit via a Turnbull NPMLE that serves as a robustness check on the baseline all-bidders estimator. Two items within the same CADMAT class can differ in scale in ways the reference price does not fully absorb, so I separate out an auction-level unobserved heterogeneity component $a_t$ following \citet{krasnokutskaya2011}'s shrinkage procedure; the resulting cost primitives are UH-clean. Entry is handled in the \citet{atheyseira2011} spirit: observed pre- and post-period bidder counts per auction are taken as equilibrium arrival rates, and the main specification allows the post-policy SME bidder distribution to differ from the pre-policy one as part of the equilibrium-selection object.

Three findings emerge.

\noindent\textbf{(i) The structural decomposition is stable across classes: the intensive margin dominates.} Let $p_{S_1}$ denote the simulated winning price under the pre-period open regime, $p_{S_2}$ the SME-only counterfactual holding the SME pool at its pre-period size, and $p_{S_3}$ the SME-only equilibrium with endogenous entry. In the main specification, $p_{S_3} - p_{S_1} = +0.24$ of the reference price in non-pharma and $+0.34$ in pharma, and the intensive shares are 77 and 75 percent. The Turnbull NPMLE gives 74 and 81 percent; the strict-invariance benchmark in the appendix gives 85 and 79 percent. Across product classes and identifying choices, 75--85 percent of the effect is intensive. The set-aside primarily changes the second-order statistic inside the auction, not the participation margin.

\noindent\textbf{(ii) Entry is the government's partial insurance.} The extensive margin matters, but mainly because it dampens what would otherwise be a larger price shock. In the counterfactual that imposes the full set-aside while holding the SME pool fixed at its pre-policy size, the price effect is 162 percent of the realized effect in non-pharma and 170 percent in pharma. Without endogenous entry, the set-aside's welfare cost would be roughly 60--70 percent larger. The policy is therefore partially self-correcting: the SMEs induced to enter by the protected segment absorb part of the intensive shock, though far from all of it.

\noindent\textbf{(iii) The welfare ranking of set-asides versus preferences is goods-characteristic-dependent.} A 10 percent price preference is robustly superior in thicker, more standardized non-pharma markets: under the main specification it raises prices by only +0.015 relative to the open regime, versus +0.245 under the set-aside, and under strict invariance the ranking remains the same. In thinner, more heterogeneous pharma markets, the preference is Pareto-superior under the main equilibrium-selection reading but the ranking reverses under strict invariance, where the implied welfare weight needed to justify the set-aside falls to 0.7. I read that bifurcation as informative about the mechanism, not as fragility: bidder-exclusion rules are most model-sensitive when the protected pool is thin and the good is heterogeneous.

The paper makes three contributions in the same order. First, it delivers a within-auction structural decomposition of an SME preference in a developing-country procurement market, closing a gap that dates to \citet{marion2007} and has survived the subsequent auction-preferences literature (\citealp{athey2013}; \citealp{nakabayashi2013}) largely because price-only data conflate entry with bidding by construction. Second, it shows that endogenous entry is not a side issue but part of the welfare object: the entry response is the government's partial insurance, and fixed-pool counterfactuals misstate the realized policy cost. Third, it turns the policy comparison into a market-design result rather than a universal ranking. The 10 percent price preference dominates in thick, standardized markets; in thin, heterogeneous markets the ranking depends on how entry selection is modeled. That is a higher-order implication about bidder exclusion, not just about pharmaceuticals.

\paragraph{Structural lineage.} The project has three versions and this is the third. v1, the reduced-form predecessor, established the 10--11 percent price effect by staggered difference-in-differences against the 76 never-treated product groups; it delivered the mechanism question this paper resolves. v2 was a structural pilot that identified $F_c^k$ from losers-only drop-outs and used a fixed-pool counterfactual; it established the methodological feasibility but left two biases in the welfare number (the losers-only tail and the absent tax-distortion term). v3 corrects both biases, adds the endogenous-entry counterfactual, adds the Turnbull robustness, and makes the equilibrium-selection versus strict-invariance distinction explicit. Readers familiar with v1 will find the reduced-form results unchanged; readers familiar with v2 will find the structural estimates smaller in the main specification than under strict invariance precisely because entry selection is now treated as part of the equilibrium object.

A practical example fixes the magnitudes. On February 6, 2018, one month before the policy change, a S\~ao Paulo public buyer procured fifty thousand units of a CMED-regulated generic (furosemide 40 mg, item code 110639) through an open Preg\~ao. Ten firms bid, the final price sat roughly twelve percent below the reference. Eight months into the post-period, a comparable purchase by the same buyer attracted three SMEs, and the winning price cleared just above reference. The gap is the 34 percentage points the model attributes on average to the pharma case under the main specification. Most of it is the intensive shift created by bidder exclusion; part of it is offset by the new SMEs the rule induces to show up. The welfare math in Section~\ref{sec:welfare} prices what that difference costs and when it is worth bearing.

The rest of the paper proceeds as follows. Section~\ref{sec:institutional} describes the institutional setting. Section~\ref{sec:data} describes the data and sample. Section~\ref{sec:model} specifies the model. Section~\ref{sec:identification} develops identification and estimation. Section~\ref{sec:results} reports the main decomposition. Section~\ref{sec:welfare} reports the welfare decomposition. Section~\ref{sec:robustness} reports robustness. Section~\ref{sec:discussion} discusses policy. Section~\ref{sec:conclusion} concludes.
